% !TEX program = xelatex  (compile with:  tectonic resume_template.tex)
% ============================================================================
%  脱敏简历模板 · 1:1 复刻原版式 · 仅替换必要 PII(姓名/电话/邮箱/微信/GitHub号/学校/个人域名)
%  技术内容与写法均按原样保留。要填成自己的信息,改下面这几个 \newcommand 即可。
% ============================================================================
\documentclass[10pt]{article}

% ---------- 可填写的脱敏占位符(改这里即可) ----------
\newcommand{\CName}{苏亦宁}                          % 虚构姓名(占位,自行替换)
\newcommand{\CPhone}{138-0000-0000}
\newcommand{\CEmail}{your.name@example.com}
\newcommand{\CWeChat}{your\_wechat}
\newcommand{\CGitHub}{your-github}                 % 个人 GitHub 用户名
\newcommand{\CSchool}{江海大学 · 软件工程专业}        % 虚构校名+专业(占位,自行替换)
\newcommand{\CProdA}{your-product-a.com}            % 个人产品域名
\newcommand{\CProdB}{your-product-b.com}
% ------------------------------------------------------

\usepackage[a4paper,left=1.4cm,right=1.4cm,top=1.1cm,bottom=1.1cm]{geometry}
\usepackage{ctex}
\usepackage{fontspec}
\usepackage{xcolor}
\usepackage{enumitem}
\usepackage{pifont}
\usepackage{graphicx}
\usepackage[most]{tcolorbox}
\usepackage{hyperref}
\usepackage{amssymb}

% ---- 字体 ----
% 原版观感:西文 = Times New Roman,中文 = 宋体,等宽 = Consolas。
% 若本机装有原版字体,取消下面 4 行注释即可"精确复刻"(否则用下方免费等价体):
% \setmainfont{Times New Roman}
% \setCJKmainfont{SimSun}                 % Windows 宋体;macOS 改 Songti SC;跨平台可用 Source Han Serif SC
% \setmonofont{Consolas}                  % 无则 Menlo(macOS)/ Cascadia Mono
% \setCJKsansfont{Microsoft YaHei}        % 黑体系(可选,用于小节标题)
%
% 默认:跨平台可编译的免费等价体(观感接近原版,tectonic/多数发行版直接能用)
% Scale 整体缩放字号以收进 2 页(想更大/更小改这个数即可)
\newcommand{\FScale}{0.80}
\setmainfont{Liberation Serif}[Scale=\FScale]   % 与 Times 度量兼容
\setsansfont{Liberation Sans}[Scale=\FScale]
\setmonofont{DejaVu Sans Mono}[Scale=0.80]
\setCJKmainfont{Noto Serif CJK SC}[Scale=\FScale]  % 思源宋体(宋体风格)
\setCJKsansfont{Noto Serif CJK SC}[Scale=\FScale]

% ---- 配色 ----
\definecolor{accent}{HTML}{1F5FBF}   % 主蓝
\definecolor{barbg}{HTML}{EEF0F2}    % 项目头灰条
\definecolor{barpink}{HTML}{FBECEC}  % 电商粉条
\definecolor{codegray}{HTML}{F2F2F2}
\definecolor{linktxt}{HTML}{1F5FBF}

\hypersetup{colorlinks=true, urlcolor=linktxt, linkcolor=linktxt}
\setlength{\parindent}{0pt}
\setlist{nosep, leftmargin=1.4em, topsep=1pt, itemsep=1pt}
\linespread{0.96}

% ---- 内联等宽 ----
\newcommand{\code}[1]{\colorbox{codegray}{\small\ttfamily #1}}
% ---- 蓝色小节标题 + 横线 ----
\newcommand{\rsection}[1]{\vspace{9pt}{\color{accent}\sffamily\bfseries\large #1}\par\nobreak\nointerlineskip\vspace{2.5pt}\noindent{\color{accent}\rule{\linewidth}{1.6pt}}\par\vspace{4pt}}
% ---- 蓝色项目标题(带下划线) ----
\newcommand{\ptitle}[1]{\vspace{2pt}\par{\color{accent}\sffamily\bfseries #1}\par\vspace{1pt}}
% ---- 项目头色块:纯色填充、无外边框、微圆角 ----
\newtcbox{\barbox}[1][barbg]{on line, boxrule=0pt, frame hidden, colback=#1, colframe=#1, arc=2pt, outer arc=2pt, left=6pt,right=6pt,top=1.6pt,bottom=1.6pt, width=\linewidth, halign=flush left, before skip=0pt, after skip=0pt}
\newcommand{\cbar}[4][barbg]{\vspace{4pt}\barbox[#1]{#2\; \textbf{#3}\hfill {\small\color{gray!140} #4}}\vspace{2pt}\par}
% ---- 真 logo(素材库在 assets/logos/*.pdf,矢量) ----
\graphicspath{{assets/logos/}}
\newcommand{\logo}[1]{\raisebox{-0.28em}{\includegraphics[height=1.05em]{#1.pdf}}}
\newcommand{\logobig}[1]{\raisebox{-0.32em}{\includegraphics[height=1.25em]{#1.pdf}}}

\begin{document}

% ============================== 抬头 ==============================
{\Huge\sffamily\bfseries \CName}\\[3pt]
\ding{37}~\CPhone \quad|\quad \ding{41}~\CEmail \quad|\quad \logo{wechat}~\CWeChat\\[2pt]
\ding{110}~\textbf{20 years old · Full-Stack AI Agent Developer · Open to opportunities}\\[2pt]
\logo{github}~github.com/\CGitHub~（\textbf{Apache Fory Committer | Maintainer of ByteDance/DeerFlow}）\\[2pt]
\ding{115}~\CSchool ~|~ 2023年-2026年（本科跳级毕业）

% ============================== 实习经历 ==============================
\rsection{实习经历}

{\color{accent}实习期间多篇论文解析及技术文章产出到内部论坛，多次登上公司《LLM技术周刊》与部门团队号优秀征文榜单，荣登 ByteTech 首页文章周榜 13 次，ByteTech 作者等级 LV8，2025 年字节跳动优秀技术创作者榜第 13 名。}
\vspace{3pt}

\textbf{\href{https://github.com/bytedance/deer-flow}{github.com/bytedance/deer-flow}} ~$\bigstar$~ stars 76k · forks 10.3k \hfill {\small\color{gray!140}\#1 Repository Of The Day}
\begin{itemize}
\item \textbf{项目简介：} 字节排名第一的开源 Long-Horizon Agent Harness 项目，能够研究、编码与创作，内置 AIOSandbox 沙箱化执行环境、可插拔 Skill 体系、跨会话长期记忆、动态 Sub-agent 调度与系统性 Context Engineering，可处理多层级 long-horizon 任务。
\item \textbf{我的职责：} 作为项目 owner 核心作者之一，主导 1.0 到 2.0 的架构升级与迭代
  \begin{itemize}[label=$\circ$]
  \item \textbf{1.0（Multi-Agent + ReAct 风格子 Agent）：} Coordinator（意图识别）$\rightarrow$ Planner $\rightarrow$ Research Team（Researcher / Coder 支持 MCP 动态扩展外部工具）$\rightarrow$ Reporter（汇总上下文生成最终报告）；Prompt 由 OpenAI Meta Prompt 生成。
  \item \textbf{2.0（Agent Teams + Middleware Chain + Skill System + Context Engineering）：}
    \begin{itemize}[label={\raisebox{0.15ex}{\scriptsize$\blacksquare$}}]
    \item \textbf{Agent Teams Harness：} 相比 1.0 五节点固定流水线，2.0 改为\textbf{单一 Lead Agent 统一决策}；\code{system\_prompt} \textbf{动态写入}——按当前启用的 Skill 列表、跨会话持久化记忆、Sub-agent 并发规则组装；Lead Agent 通过 \code{task} 工具触发 SubagentExecutor 新起独立 Agent 实例；\textbf{调度与执行双线程池分离}，最多 \textbf{3 并发 / 单任务超时 900s}；Sub-agent 继承父线程目录 + 复用父沙箱，每步\textbf{实时推流至前端}，完成后压缩摘要回传 Lead Agent。
    \item \textbf{Middleware Chain：} 1.0 上下文压缩 / 沙箱 / 记忆逻辑散落各处，改一处要变动多处；2.0 抽成 \textbf{11 层有序流水线}，新增或修改能力只对应层，Agent 推理代码零改动（ThreadData $\rightarrow$ Uploads $\rightarrow$ Sandbox $\rightarrow$ DanglingToolCall $\rightarrow$ Summarization $\rightarrow$ Todo $\rightarrow$ Title $\rightarrow$ Memory $\rightarrow$ ViewImage $\rightarrow$ SubagentLimit $\rightarrow$ Clarification）。
    \item \textbf{Skill System：} 每个 Skill 是一个目录，内含 \code{SKILL.md}；\code{system\_prompt} 只注入索引，Agent 按需 \code{read\_file} 渐进加载，避免全量注入 token 爆炸；支持自定义 Skill，\code{custom/} 下新增目录即可扩展。
    \item \textbf{Context Engineering：} \textbf{Write}——中间结果写入文件系统，索引 / 规则每轮注入 \code{system\_prompt}；\textbf{Select}——按 token 预算从 \code{memory.json} 裁剪注入记忆，Skill 按需读全文；\textbf{Compress}——超限调轻量 LLM 压缩历史 + 异步提炼持久化记忆；\textbf{Isolate}——Sub-agent 仅见任务描述，结果摘要回传，沙箱按 \code{thread\_id} 隔离。
    \end{itemize}
  \end{itemize}
\end{itemize}

% ---------- 抖音 AI 基建 ----------
\cbar{\logo{douyin}}{抖音 AI 基建 - AI4SE 基座算法 \& Infra}{2025.12 - 2026.4}

\ptitle{AgentScaling for SFT \& RL~|~MultiModel Relay + Multi-Scaffold Diverse Trajectory Synthesis Pipeline}
\begin{itemize}
\item \textbf{背景与目标：}
  \begin{itemize}[label=$\circ$]
  \item 单一模型轨迹风格单一易过拟合——借鉴「Mix Diverse GPTs」，引入多种 scaffold 提供风格不同的轨迹，artifacts 生成任务上混合训练已验证有显著收益。
  \item 针对 Claude Opus 4.6 连续失败 3 次以上的难题，用接力和定位策略提升通过率、降低重复推理成本。\textbf{目标：提升难题通过率，降低重复推理成本}。
  \end{itemize}
\item \textbf{我的职责：项目 owner，完全负责从 0 到 1 设计并落地以下系统架构与技术链路}（CLI: \code{relay run $\rightarrow$ export}）：
  \begin{itemize}[label=$\circ$]
  \item \textbf{多 scaffold 统一适配：} 统一 Agent 接口接入 10+ scaffold（codex / claude\_code / cline 等），新增 scaffold 实现统一 run / get\_answer 接口即可接入。
  \item \textbf{relay run 执行引擎：} 同一 Docker 容器内多 LLM 顺序接力，共享代码现场。切换策略为\textbf{三阶段交接}：规则压缩历史为 \textasciitilde12k 结构化文本 $\rightarrow$ LLM 以 gold patch 为隐形顾问反出 RFC（复盘）+ Plan A/B/C $\rightarrow$ 裁剪对话为 [首条任务描述 + handoff 消息] 作为下一模型输入；可选每轮 \textbf{Test-in-Loop} 提取 patch 评测，通过则提前退出，失败则注入失败报告给下一模型。续跑时从 checkpoint 恢复仓库 patch 与对话，从上一模型继续。针对反复失败的难题，复现 SWE-Replay 论文并提升策略：维护历史轨迹 Archive，从关键中间步骤分叉、apply git diff 恢复现场后重试，避免全量重跑。\textbf{Benchmark 指标：} SWE-Bench 系列难题通过率较全量重跑提升 25\%+，单任务推理成本（API 与时长）降低约 30\%+，多 scaffold 混合轨迹覆盖 10+ 种风格。
  \item \textbf{上下文压缩 + Checkpoint 续跑：} 超 180k 触发 Micro compact，超 250k 触发 LLM 兜底；每轮 checkpoint，中断后可从上次模型续跑。
  \end{itemize}
\end{itemize}

\ptitle{Agent Tracer Evaluation for SFT \& RL~|~Trajectory Diagnosis + Rubric Eval and Wash Pipeline}
\begin{itemize}
\item \textbf{背景与目标：}
  \begin{itemize}[label=$\circ$]
  \item 多模型 MultiSWEBench 完成率接近，但 Agentic 能力（规划 / 工具链效率 / 错误恢复 / patch 质量）参差不齐；\textbf{目标：用 Rubric 评判标准逐 check 打分，筛选高质量轨迹片段用于 SFT / RL 训练}。
  \end{itemize}
\item \textbf{我的职责：项目 owner，完全负责从 0 到 1 设计并落地以下系统架构与技术链路}（CLI: \code{evaluate run $\rightarrow$ wash $\rightarrow$ export}）：
  \begin{itemize}[label=$\circ$]
  \item \textbf{LLM-as-Judge 评估引擎：} 按 check 类型分组 batch 减少 API 调用，规则类 check 零 LLM 直接判断，失败时回退逐条补全；支持多模型投票；遇 429 自动退避重试，并发数与请求频率双重限流。
  \item \textbf{Failure Onset 定位：} Agent 失败根因是能定位相关代码证据但无法正确解读并转化为行动（\textbf{Evidence-to-Action Gap}）；通过将 Rubric check 逐步骤打分定位到具体轮次，自动识别 failure onset（首次出现错误决策的步骤），过滤 onset 后的级联失败片段，只保留 onset 前的高质量决策轨迹用于训练。
  \item \textbf{轨迹压缩：} 逐轮按内容类型差异化截断；超 180k token 触发 \code{micro\_compact}；超 250k 触发 LLM 总结兜底，保留首条 USER 消息与失败报告。
  \item \textbf{增量续跑 + wash / export：} 自动增量补全未完成 check；按 CSR / ISR / Score / Effective Action Ratio（有效工具调用占比）多维 Top-K 筛选 + 早期质量评估，导出 pkl 格式。\textbf{Benchmark 指标：} Agentic Rubric 模型区分度 35\%，Failure Onset 定位后高质量决策片段占比提升，用于 SFT/RL 的轨迹通过率与 patch 质量显著优于未筛选基线。
  \end{itemize}
\end{itemize}

\ptitle{SWE 等系列 Benchmark 评测服务化改造~|~接入字节云 AIPaaS 容器平台}
\begin{itemize}
\item \textbf{我的职责：} 项目 owner，将分散的本地评测接入 AIPaaS 云容器统一入口：spec 配置 $\rightarrow$ \code{CreateSession} 申请隔离容器 $\rightarrow$ 还原仓库 $\rightarrow$ Agent 修复 $\rightarrow$ diff 送测 $\rightarrow$ 异步轮询 $\rightarrow$ 返回结构化结果，供各业务按需调用。
\end{itemize}

% ---------- 抖音电商 ----------
\cbar[barpink]{\logo{douyin}}{抖音电商 - 营销 \& 互动玩法}{2025.5 - 2025.9 ~|~ 产运 Agent~|~ B 端研发平台 \& C 端互动玩法频道页 ~\href{https://github.com/\CGitHub/ecom}{→github.com/\CGitHub/ecom}}

\ptitle{EC-Monica 电商产运 Agent}
\begin{itemize}
\item \textbf{背景：} 为产品、运营、研发提供电商业务 Agent，支持数据分析与报告生成、业务归因与需求文档解析等任务，提高业务理解效率与文档编写效率。
\item \textbf{指标与效果：} 上线后服务产运、产品、研发等多角色，调用量月均 600+ 次，报告类任务占比约 55\%，需求 / 归因解析类约 45\%；单次分析报告生成从人工 1--2 小时缩短至 5--15 分钟；业务范文 Few-shot 注入后，输出风格与内部模板对齐率提升，复购与主动推荐率持续增长。
\item \textbf{我的职责：项目 owner，完全负责产品从 0 到 1 设计并落地以下内容}
  \begin{itemize}[label=$\circ$]
  \item \textbf{Agentic Workflow：} 意图识别 $\rightarrow$ 条件路由（需工具 / 直接生成）$\rightarrow$ \code{cfg.json} 动态注册 FunctionSchema 的 Tool Registry 选工具循环执行 $\rightarrow$ 数据预处理 / 计算同比 / 环比指标后注入 Context $\rightarrow$ LLM 生成子结论 $\rightarrow$ 循环直至完成后汇总报告；工具层封装超时重试与错误分类，失败时 Fallback 直接生成。
  \item \textbf{Context Engineering：} 注入业务范文与运营分析模板作为 Few-shot 样例强化风格对齐；多轮推理中对 Observation 做摘要压缩 + 滑动窗口截断，携带压缩后的前轮结论滚动输入。
  \end{itemize}
\end{itemize}

% ---------- 抖音直播中台 ----------
\cbar{\logo{douyin}}{抖音直播中台}{2025.1 - 2025.5 ~|~ AI 智能搭建~|~ C 端直播玩法开发 \& B 端平台权限体系 ~\href{https://github.com/\CGitHub/live_stream}{→github.com/\CGitHub/live\_stream}}

\ptitle{直播 AI 智能搭建平台 — Agentic Workflow 搭建}
\begin{itemize}
\item \textbf{背景：} 基于 GenUI 平台，为直播搭建团队打造 AI 驱动的低代码智能建站能力，实现"自然语言 $\rightarrow$ 页面 / 流程"的端到端自动化搭建，显著降低页面搭建门槛与人力成本。\textbf{使用率与效果：} 上线后 AI 搭建能力在直播搭建场景中的使用率占新建 / 改版页面的约 40\%+，单页搭建耗时从人工 2--4 小时降至 15--30 分钟，团队采纳率与复购率持续提升。
\item \textbf{我的职责：项目 owner，完全负责产品从 0 到 1 设计并落地以下内容}
  \begin{itemize}[label=$\circ$]
  \item 设计并实现多条 Agentic Workflow 链路与 Prompt 工程设计：
    \begin{itemize}[label={\raisebox{0.15ex}{\scriptsize$\blacksquare$}}]
    \item \textbf{Prompt 工程：} 覆盖意图分析、任务拆分（new\_page / new\_section / update\_section / new\_flow / bind 等 8 种动作）、数据源缺失检测等完整推理链路。
    \item \textbf{主流水线：} 负责意图分析，将用户输入与页面 UIDL 传入 Claude3.5\_v2，输出结构化 taskList（含 UI / Flow / Bind 三类任务）。
    \item \textbf{Controller：} 按优先级并行调度各子任务，通过 Promise.allSettled 实现多任务并发执行；封装 XML Patch 服务实现 UIDL 差量更新。
    \item \textbf{UI 生成链路：} 支持文生 / 图生双模式，经 JSX $\rightarrow$ XML $\rightarrow$ UIDL 多级 IR 转换渲染完成。
    \item \textbf{Flow 生成链路：} 将 AI 输出的 IRFlow 转换为可执行的流程结构，完成节点与边的构建。
    \item \textbf{Bind 链路：} 负责将 UI 组件与 Flow 按事件类型（onClick / onMount 等）绑定到 UIDL flowMap。
    \end{itemize}
  \item 一期向 AI 喂入 54 个基础组件的 Schema/DSL 定义完成模型训练，为后续接入 Semi / Arco Design 组件库奠定基础。
  \end{itemize}
\end{itemize}

% ---------- B 站 ----------
\cbar{\logo{bilibili}}{哔哩哔哩 - OGV 技术部}{2024.8 - 11 ~|~ 微前端低代码流程工作台~|~ B/C 端业务开发 ~\href{https://github.com/\CGitHub/bilibili}{→github.com/\CGitHub/bilibili}}

% ============================== 兼职 & 开源 ==============================
\rsection{AI Startup 兼职 \& 开源项目 \& AI 黑客松比赛：\normalsize\href{https://github.com/\CGitHub/Awards}{github.com/\CGitHub/Awards}}

\ptitle{Thumbnail AI · Mountex}
\begin{itemize}
\item \textbf{背景：} 基于 MkSaaS、Next.js 搭建的多平台封面图生成 AI 工具（\href{https://\CProdA}{\CProdA}），面向 YouTube、Instagram、Facebook、X、LinkedIn、TikTok 等多平台内容创作者，支持选择平台尺寸、风格参考与主体上传，三步内完成专业封面图生成，帮助创作者提升点击率。
\item \textbf{产品与链路：} 前端 Next.js + 响应式多端适配，支持平台预设尺寸一键选择、本地上传 / URL 主体、风格参考图；AI 链路：主体与风格图输入 $\rightarrow$ 提示词编排（平台场景 + 风格关键词）$\rightarrow$ 调用图像生成 API 生成多尺寸封面 $\rightarrow$ 预览与下载；支持免费额度 + 订阅付费，付费用户可解锁高清导出与批量生成。
\item \textbf{数据与指标：} 累计注册用户 2000+，付费用户 200+；月活 500+，月生成封面 1500+ 张；付费转化率约 38\%；用户从上传到出图平均时长 < 2 分钟。
\item \textbf{我的职责：} 项目 owner，从 0 到 1 负责前端工程（Next.js、支付与订阅流程、多平台尺寸与风格配置）、AI 封面生成链路设计与对接（提示词工程、多尺寸批量生成、降本与体验优化）。
\end{itemize}

\ptitle{Krene-Art}
\begin{itemize}
\item \textbf{背景：} 面向游戏 / 二次元角色设计师的 AI 创作平台（\href{https://\CProdB}{\CProdB}），通过 Agent 完成意图解析、参考图检索、视觉特征提取与 Prompt 生成，输出带设计说明的角色卡片（含核心设定、概念整合与角色立绘），缩短从脑洞到成图的设计周期。
\item \textbf{ReAct Workflow 与链路：}
  \begin{itemize}[label=$\circ$]
  \item \textbf{search：} 根据用户意图生成检索关键词 $\rightarrow$ 爬取 Bing 图搜解析高清链接 $\rightarrow$ LLM 按游戏角色维度（造型、配色、气质等）打分筛选，单次聚合约 300 张参考。
  \item \textbf{design：} GPT-4o Vision 对参考图做五维视觉特征提取 $\rightarrow$ 与用户设定拼接成结构化 Prompt $\rightarrow$ NovelAI nai-diffusion-3 生成约 10 张角色图。
  \item \textbf{evaluate：} 按设计师设定维度打分，低分时 ReAct 触发重新 design（调整 Prompt 或参考集），直至达标后输出结构化角色卡片。
  \end{itemize}
\item \textbf{数据与指标：} 累计用户 1200+，付费用户 150+；月生成角色卡 400+ 张；付费转化率约 42\%；单角色从输入意图到出卡平均 6--8 分钟。
\item \textbf{我的职责：} 项目 owner，从 0 到 1 负责前端工程与 AI 角色设计 ReAct Workflow（意图解析、search/design/evaluate 三阶段编排、Vision 特征提取与 NovelAI 对接、角色卡片结构化输出与展示）。
\end{itemize}

\begin{itemize}
\item \textbf{开源：} 持续半年 Apache 社区贡献，\textbf{18 岁成为全球最年轻 Apache Committer}；深度参与 ByteDance/DeerFlow、Apache Fory、Spring AI Alibaba、Ant Design X 等社区开发与维护；2024 阿里天池云原生全球编程挑战赛 issue 解决数第一，技术文章荣获天池赛最佳质量奖；2024 腾讯犀牛鸟开源大赛课题实战奖 \& Issue 实战奖；2024 OSPP 中科院开源之夏（ByteDance/VisActor 组件开发）。
\item \textbf{AdventureX AI 黑客松：} 开发全球首个 Apple Vision Pro Agent，斩获 Kimi For Vibe Coding 奖、空间智能赛道第二名及 Injective \$1000。
\item \textbf{AI Startup 好朋友：} 帮多家苗格、红杉、蓝驰、奇绩等知名 VC 被投 AI Startup 实现 Agent 业务落地；与千问、Kimi、智谱 GLM、MiniMax、豆包、Trae、灵光、阿福等多家大模型 / AI 产品深度合作，参与官方评测与 Demo 制作，产出技术文章、评测笔记与 Demo 视频，为产品宣传与迭代提供反馈。
\item 高中时期做过 Live2D/3D 数字人，二次元平台产品经理，对二次元文化有深入了解。
\end{itemize}

\end{document}
